\section{Sixth Wound: The Saturated Shadow}

The five wounds above were handed to me. Most are the historical objections, raised by people who wanted the mechanical program to succeed, and every one of them is older than I am. This one nobody handed me: I found it while writing the section that closes this chapter, and it is the least settled thing in the book. So I will set it out one step at a time.

Begin with something the shadow push has needed since its first page and has so far been allowed to have without comment. Gravity, here, is a push from a rain that is blocked: two bodies approach because each removes some of the rain that would otherwise have struck the other from that side. Now ask what the size of that push depends on. If a body stops only a very small fraction of what passes through it, then every piece of matter in it does its own blocking independently, nothing hides behind anything else, and the total shadow is the sum over all the matter. The push then tracks the body's \emph{mass}. But if a body stops a large fraction, the far side of it stands in the near side's shadow and contributes nothing. The body is then blocking as much as it can, and the push tracks the \emph{silhouette} --- the area of the outline --- and stops caring how much matter is stacked up behind it.

Two regimes, then, and everything turns on which one the world is in. Physicists say optically thin and optically thick; I will say transparent and opaque. \emph{The shadow push works only in the transparent regime.} That is the requirement. The rest of this section is about whether the world meets it.

Everywhere we can weigh, it does, and obligingly. The Earth is a ball of rock eight thousand miles thick and its gravity still counts every atom in it --- which is as direct a measurement as anyone could want of how little of the rain is stopped on the way through. Nor does the testing stop at ordinary densities. Binary pulsars, neutron stars orbiting one another and keeping time better than most clocks, have their masses and their orbits checked against each other to many decimal places, and the checks pass. So matter is still transparent at something like four times ten to the seventeenth kilograms in every cubic metre --- the density of an atomic nucleus. Nature has tested the requirement for us, hard, and it has held.

The question is whether it holds all the way. What of a body that has fallen past the last rung we know of, something dense enough to trap light?

Suppose such a body is opaque. Then the numbers come out wrong, and it is worth being careful about \emph{how} wrong, because the first version of this section was wrong in the toy's disfavour and I would rather be accurate than frightening. There is a famous relation in which a black hole's radius grows in direct proportion to its mass, and I reached for it. But that relation describes a \emph{horizon}, and a horizon does no blocking; blocking is done by matter. A horizon is in fact a poor stand-in for matter, as its own numbers confess: the average density inside one falls as the inverse square of the mass, so that a hole of ten suns encloses matter at roughly nuclear density, the hole at this galaxy's centre encloses an average nearer that of a heavy oil, and a hole of ten thousand million suns encloses an average thinner than the air in this room. Nothing thinner than air holds anything up. Whatever sits inside a large hole must sit far inside it, on a strut of its own, at a density set by that strut and not by the horizon.

Take the body seriously and the arithmetic is the ordinary arithmetic of solid things. We live in three dimensions, mass fills volume, and a body held at some characteristic stiffness has a radius growing as the cube root of its mass and therefore a silhouette growing as about the two-thirds power. Possibly weaker: real degenerate bodies are perverse about this, and white dwarfs actually \emph{shrink} as they gain mass. So an opaque collapsed body would push as roughly the two-thirds power of its mass, where the world insists the push grows as the mass itself. The two rules diverge as the cube root of the mass ratio --- a factor of about seventy-five between a hole of ten suns and the one at the centre of this galaxy, and about a thousand across the whole observed range of holes.

That is a discrepancy the sky would have caught. The mass of the dark object at our galaxy's centre can be read from the orbits of the stars swinging around it, and independently from the size of the ring the Event Horizon Telescope photographed; the waveforms of inspiralling pairs match general relativity using one mass throughout. A push that tracked outline rather than mass would show up as those numbers refusing to agree with one another. They agree.

So here is the wound, as plainly as I can put it. The shadow push requires matter to be at least nearly transparent to gravity's massions. That requirement is verified up to nuclear density and not one step beyond. At the only place beyond it that we can observe, the observations say the requirement still holds --- and the toy has no right to expect that and cannot explain it.

The best reply I have is real, and it is not free. A horizon is our layer's \emph{optical} fact: it is defined by what our light cannot climb out of. The massions are not our light. They belong to a layer below ours, they move very much faster, and they are under no obligation whatever to be stopped by whatever stops a photon. A collapsed body may therefore be black and still nearly transparent --- opaque to us, and hardly there at all as far as the rain is concerned. Then the push tracks the mass, the orbits come out right, and the wound closes.

It closes, though, by assuming exactly what is in question. Nothing here \emph{derives} the transparency; it is asserted, and asserted about the one regime in which nothing can be measured. Worse, the assertion is a demanding one. The massion's chance of being stopped, per unit of matter it passes through, must be small enough that even matter crushed past the neutron rung stays optically thin --- and it must be small for a reason, not because the toy requires it. Every wound in this chapter has cost the toy one hidden thing. This one costs it a hidden \emph{number}, and hidden numbers chosen to fit are how mechanical theories have historically died.

What would settle it is that number, measured or derived: the massion's stopping chance per unit of matter, and with it the density at which the shadow saturates. If saturation lies far beyond even collapsed matter, the toy keeps its proportionality to mass and owes an explanation of why the number is so small. If it does not, then the gravity of collapsed bodies must somehow be made to track their mass rather than their outline, and I do not see how. That is the sixth problem of the research program below, and of the six it is the one I would most like to see somebody take away from me.

